% GENERATED FILE. Edit canonical lessons, then run npm run generate:books.
% canonical-source-hash: fnv1a64:40157ff5ea1dce74
% canonical-lessons: JA-C78-doko, JA-C78-chikaku, JA-C78-ushiro, JA-C78-yoko, JA-C78-mannaka

\chapter{Where, Nearby, Behind, Beside, Very middle}
\label{ch:ja-where-nearby-behind-beside-very-middle}
\hlchaptermodality{78}

\begin{chapteropening}
\textbf{By the end of this chapter:} \emph{I can say where, nearby, behind, beside and the very middle.}

Complete the last lesson of chapter 78: I can say where, nearby, behind, beside and the very middle.
\end{chapteropening}

\section[\texorpdfstring{\ja{どこ} (doko)}{どこ (doko)}]{\ja{どこ} (doko) — where}

\label{lesson:JA-C78-doko}



\begin{quote}
Before the new one: say the Japanese for the fifth of the month; some day, then the Japanese for the sixth of the month.
\end{quote}

\subsection*{You'll want to know: \ja{どこ}}
\textbf{\ja{どこ}} — \emph{doko} — \textquotedblleft{}where\textquotedblright{}.

\begin{grammarlens}[title={What you've built}]
One more word for this chapter.
\end{grammarlens}

\subsection*{Your turn}
\begin{itemize}
\raggedright
  \item \emph{Say it:} \emph{doko}
  \item \emph{Say it:} \emph{doko}, once more
  \item \emph{Say it:} say \emph{doko}
  \item \emph{Recall:} say the Japanese for the fifth of the month; some day, then the Japanese for the sixth of the month, then say \emph{doko} again
\end{itemize}

\begin{tcolorbox}[breakable,colback=teal!4,colframe=teal!35!black,title={Before you move on}]
What does \ja{どこ} mean? (\textquotedblleft{}Where\textquotedblright{}.) Say it once more.
\end{tcolorbox}
\section[\texorpdfstring{\ja{ちかく} (chikaku)}{ちかく (chikaku)}]{\ja{ちかく} (chikaku) — nearby}

\label{lesson:JA-C78-chikaku}



\begin{quote}
Before the new one: say the Japanese for the sixth of the month, then the Japanese for where.
\end{quote}

\subsection*{You'll want to know: \ja{ちかく}}
\textbf{\ja{ちかく}} — \emph{chikaku} — \textquotedblleft{}nearby\textquotedblright{}.

\begin{grammarlens}[title={What you've built}]
One more word for this chapter.
\end{grammarlens}

\subsection*{Your turn}
\begin{itemize}
\raggedright
  \item \emph{Say it:} \emph{chikaku}
  \item \emph{Say it:} \emph{chikaku}, once more
  \item \emph{Say it:} say \emph{chikaku}
  \item \emph{Recall:} say the Japanese for the sixth of the month, then the Japanese for where, then say \emph{chikaku} again
\end{itemize}

\begin{tcolorbox}[breakable,colback=teal!4,colframe=teal!35!black,title={Before you move on}]
What does \ja{ちかく} mean? (\textquotedblleft{}Nearby\textquotedblright{}.) Say it once more.
\end{tcolorbox}
\section[\texorpdfstring{\ja{うしろ} (ushiro)}{うしろ (ushiro)}]{\ja{うしろ} (ushiro) — behind}

\label{lesson:JA-C78-ushiro}



\begin{quote}
Before the new one: say the Japanese for where, then the Japanese for nearby.
\end{quote}

\subsection*{You'll want to know: \ja{うしろ}}
\textbf{\ja{うしろ}} — \emph{ushiro} — \textquotedblleft{}behind\textquotedblright{}.

\begin{grammarlens}[title={What you've built}]
One more word for this chapter.
\end{grammarlens}

\subsection*{Your turn}
\begin{itemize}
\raggedright
  \item \emph{Say it:} \emph{ushiro}
  \item \emph{Say it:} \emph{ushiro}, once more
  \item \emph{Say it:} say \emph{ushiro}
  \item \emph{Recall:} say the Japanese for where, then the Japanese for nearby, then say \emph{ushiro} again
\end{itemize}

\begin{tcolorbox}[breakable,colback=teal!4,colframe=teal!35!black,title={Before you move on}]
What does \ja{うしろ} mean? (\textquotedblleft{}Behind\textquotedblright{}.) Say it once more.
\end{tcolorbox}
\section[\texorpdfstring{\ja{よこ} (yoko)}{よこ (yoko)}]{\ja{よこ} (yoko) — beside}

\label{lesson:JA-C78-yoko}



\begin{quote}
Before the new one: say the Japanese for nearby, then the Japanese for behind.
\end{quote}

\subsection*{You'll want to know: \ja{よこ}}
\textbf{\ja{よこ}} — \emph{yoko} — \textquotedblleft{}beside\textquotedblright{}.

\begin{grammarlens}[title={What you've built}]
One more word for this chapter.
\end{grammarlens}

\subsection*{Your turn}
\begin{itemize}
\raggedright
  \item \emph{Say it:} \emph{yoko}
  \item \emph{Say it:} \emph{yoko}, once more
  \item \emph{Say it:} say \emph{yoko}
  \item \emph{Recall:} say the Japanese for nearby, then the Japanese for behind, then say \emph{yoko} again
\end{itemize}

\begin{tcolorbox}[breakable,colback=teal!4,colframe=teal!35!black,title={Before you move on}]
What does \ja{よこ} mean? (\textquotedblleft{}Beside\textquotedblright{}.) Say it once more.
\end{tcolorbox}
\section[\texorpdfstring{\ja{まんなか} (mannaka)}{まんなか (mannaka)}]{\ja{まんなか} (mannaka) — the very middle}

\label{lesson:JA-C78-mannaka}



\begin{quote}
Before the new one: say the Japanese for behind, then the Japanese for beside.
\end{quote}

\subsection*{You'll want to know: \ja{まんなか}}
\textbf{\ja{まんなか}} — \emph{mannaka} — \textquotedblleft{}the very middle\textquotedblright{}.

\begin{grammarlens}[title={What you've built}]
One more word for this chapter.
\end{grammarlens}

\subsection*{Your turn}
\begin{itemize}
\raggedright
  \item \emph{Say it:} \emph{mannaka}
  \item \emph{Say it:} \emph{mannaka}, once more
  \item \emph{Say it:} say \emph{mannaka}
  \item \emph{Recall:} say the Japanese for behind, then the Japanese for beside, then say \emph{mannaka} again
\end{itemize}

\begin{tcolorbox}[breakable,colback=teal!4,colframe=teal!35!black,title={Before you move on}]
What does \ja{まんなか} mean? (\textquotedblleft{}The very middle\textquotedblright{}.) Say it once more.
\end{tcolorbox}
